\chapter{Introducción}\label{cap:introduccion}

\section{Motivación}

En diciembre de 2025 KEXP grabó a Angine de Poitrine en el festival Trans Musicales de Rennes; el video, publicado en
febrero de 2026, pasaba de 7.5 millones de reproducciones en YouTube a fines de marzo \citep{kexp2026}. Lo que llama la atención al escucharla es una combinación poco frecuente: una guitarra y
un bajo con trastes de cuarto de tono, canciones que se construyen capa por capa con un looper y compases que no son los
habituales del rock. Esa combinación plantea la pregunta que
guía este trabajo:

\begin{quote}
\itshape ¿La música de Angine de Poitrine encaja en un estilo existente o es una propuesta nueva?
\end{quote}

Responderla de oído es difícil, porque la banda mezcla rasgos de tradiciones muy distintas. Responderla por computador
también lo es, porque las herramientas estándar de análisis musical piensan en doce notas por octava: \texttt{music21}
comprime toda altura a doce clases y los transcriptores automáticos redondean cada cuarto de tono al semitono más
cercano. Para esta música, medir en doce notas es borrar justamente lo que la hace distinta. Por eso el trabajo empieza
por construir una herramienta que mida en veinticuatro.

\section{Qué es un análisis computacional}

Un análisis computacional asigna a una pieza una macroestructura (resultados visuales, datos o patrones) a partir de un
programa que lee la información contenida en su registro. La obra musical es un objeto abstracto, una actividad, pero sus
registros existen en contextos concretos: una grabación, una partitura, un archivo. El programa permite reconstruir de
manera empírica, y repetible, el fundamento técnico de esos registros.

La entrada de un análisis así es siempre una representación, y ninguna representación es única. Un mismo audio en WAV y
en FLAC no contiene el mismo flujo de datos; un MP3 de baja calidad pierde información que un formato sin compresión
conserva; MusicXML y MIDI hacen explícitos los compases y las alturas, pero solo los que alguien decidió escribir. La
representación de una pieza queda indeterminada de tres maneras: por cómo se proyectan los datos, por qué datos se
seleccionan y por la información que se introduce desde fuera (o la incertidumbre sobre si existe). Por eso el diagrama
de análisis necesita revisión constante para reflejar las variaciones entre entradas, incluso cuando corresponden a la
misma pieza.

En ese sentido, el análisis computacional busca cómo una pieza establece sus propios principios explicativos. Para que sea
útil, su resultado debe tener una estructura temporal, contener menos información que la pieza que resume y cumplir un
protocolo explícito que limite los sesgos de quien analiza. En este trabajo, ese protocolo es concreto: la herramienta
nunca lee la partitura, la referencia se corrige de oído y a ciegas, el azar queda fijo y el error de la propia
herramienta se mide y se aplica a los corpus con que se compara.

\section{La banda}

Angine de Poitrine es un dúo formado en 2019 en Saguenay, Quebec. Sus integrantes usan los seudónimos Khn y Klek de
Poitrine: Khn toca la guitarra y el bajo microtonales y Klek, la batería \citep{weller2026,dunlevie2026}. Su música es rock instrumental del tipo que se
suele llamar \emph{math rock}: riffs en bucle que se superponen, en compases poco comunes. Su afinación es 24-EDO, es
decir, veinticuatro notas igualmente espaciadas por octava, el doble que la guitarra común; cada traste nuevo del
instrumento agrega un cuarto de tono (50 cents). El objeto de estudio es su álbum \emph{Vol.~II} (3 de abril de 2026)
completo, en el orden de sus pistas: Fabienk, Mata Zyklek, Sarniezz, Utzp, Yor Zarad y Angor.

El nombre de la banda es el término francés para la angina de pecho, el dolor torácico causado por la falta de
oxígeno en el corazón.

\section{Influencias declaradas y comparaciones}\label{sec:influencias}

Las influencias que se ponen a prueba provienen de tres fuentes distintas, y conviene separarlas desde el principio:

\begin{itemize}
  \item \textbf{Lo que dice la banda.} En la entrevista de \emph{Noize Magazine}, Khn cuenta que al principio tiraban más
  al rock turco y que siguen escuchando rock turco y música de Medio Oriente; que usa los cuartos de tono con un fraseo
  «más cubista, más a lo Zappa»; y que la armonía de casi todos sus temas se compara con el \emph{Überjam} de John
  Scofield o con «So What» de Miles Davis. Klek dice que su banda favorita es Gentle Giant
  \citep{noize}. En el podcast \emph{Aucard de Tours}, Khn describe una música construida sobre loops, con la
  repetición del techno en la estructura y no en el timbre, y con la idea de que la gente baile «en 5 y en 7»
  \citep{aucard2026}.
  \item \textbf{Lo que dicen los analistas.} Mathémusique (Paul Lascabettes) pone como precedente a King Gizzard and the
  Lizard Wizard, que en 2017 pasó de doce a veinticuatro notas en \emph{Flying Microtonal Banana}; la banda no lo cita
  \citep{mathemusique2026}. David Bennett afirma que en uno de sus temas los cuartos de tono se usan como en el maqam
  hijaz, «casi» un frigio dominante \citep{bennett2026}.
  \item \textbf{Una comparación propia.} Primus se eligió al estudiar su armonía. Ninguna fuente pública consultada lo
  cita como influencia de la banda; en este trabajo cumple el papel de control, porque es rock sin cuartos de tono.
\end{itemize}

Cada influencia apunta a una dimensión que se puede medir: el rock turco y el maqam a la escala; Zappa y Gentle Giant a
la melodía y la métrica; Scofield y Miles Davis a la armonía; el techno y los loops a la forma en ciclos; y la
declaración «en 5 y en 7» al metro. Así la pregunta general se descompone en preguntas que se contrastan dimensión por
dimensión.

\section{Hipótesis}

\begin{quote}
La banda puede ser típica en cada dimensión (armonía, melodía y ritmo) de un estilo distinto, pero atípica en la
combinación: una recombinación nueva de lenguajes existentes.
\end{quote}

La hipótesis nula correspondiente es que, una vez descontado el error de medición de la herramienta, el álbum no se
distingue del rock. La calibración del método con piezas reales de rock (capítulo~\ref{cap:resultados}) es lo que
permite rechazarla o no.

\section{Objetivos}

\paragraph{Objetivo general.} Decidir, a partir del audio, si la música de Angine de Poitrine cabe en un estilo
existente o constituye una propuesta nueva.

\paragraph{Objetivos específicos.}
\begin{enumerate}
  \item Transcribir desde el audio la melodía, el bajo, los compases y los acordes en 24-EDO.
  \item Medir los ciclos del looper y el uso de los cuartos de tono, con un control que estime el error de la propia
  herramienta.
  \item Evaluar la transcripción contra referencias corregidas de oído, con un protocolo a ciegas.
  \item Comparar cada dimensión con corpus de estilos y con repertorios de referencia, y combinar las comparaciones en
  un veredicto calibrado.
\end{enumerate}

\section{Aportes}

\begin{enumerate}
  \item \textbf{Una herramienta:} transcripción y análisis armónico abiertos en 24-EDO, desde el audio polifónico, guiados
  por la repetición del looper. Su salida es el tabroll.
  \item \textbf{Un método:} comparación con corpus reales de estilos por bloques de rasgos, con el error de la
  herramienta aplicado a los corpus y calibrada para no exagerar la novedad.
  \item \textbf{Un resultado:} una respuesta medida a la pregunta de si la banda es un estilo existente o una propuesta
  nueva.
\end{enumerate}

\section{Organización del documento}

El capítulo~\ref{cap:marco} presenta el marco matemático y musical: la medida de la altura en 24-EDO, las herramientas de
procesamiento de señales, la teoría armónica que se extiende y las medidas estadísticas de comparación. El
capítulo~\ref{cap:estado} resume el estado del arte y el vacío que cubre el trabajo. El capítulo~\ref{cap:datos}
describe los datos: el álbum, las referencias, los controles y los corpus. El capítulo~\ref{cap:metodologia} cuenta
cómo evolucionó el método y describe el algoritmo final, su evaluación y la comparación con estilos, y el
capítulo~\ref{cap:herramienta}, cómo está construida la herramienta y qué garantiza su reproducibilidad. El
capítulo~\ref{cap:resultados} presenta los resultados, el capítulo~\ref{cap:discusion} discute sus límites y el
capítulo~\ref{cap:conclusiones} cierra con las conclusiones y el trabajo futuro.
